% !PRINTED 85
\weberpage{85}
%<<CONT
alle nicht durch $q_h^{\varkappa-1}$ theilbar, und die Zahlen
% [class ruling, 2026-08-06, 6x zoom on b2 idx 106 = printed 85]
% The glyph is VARKAPPA, the hooked form. On that page it is legible in
% q_h^{varkappa-1} and again in running text: "nur varkappa - 1 oder 2
% (bei varkappa = 2) an Stelle von varkappa". It is the curved form with no
% straight vertical stem, which is the standard German mathematical kappa of the
% period -- not x and not u. This is a property of the typeface, so it holds
% wherever the letter appears in these volumes. \varkappa is correct.
\weq{14}{i_{1,h},\ i_{2,h},\ \dots\ i_{\nu,h}}
sind alle nicht durch $q_h$ theilbar, folglich ist nach (11) keine der
Zahlen $i_1, i_2, \dots i_\nu$, und also auch $J$ nicht durch $q_h^{\varkappa-1}$ theilbar,
was den Annahmen a) und b) widerspricht. Für $q_h = 2$ ist in
diesem Schlusse nur $\varkappa - 1$ oder $2$ (bei $\varkappa = 2$) an Stelle von $\varkappa$
zu setzen, um denselben Widerspruch gegen c) und d) zu erhalten.

    Ist aber $\varkappa = 1$ und sind die Zahlen (12) alle durch $q_h - 1$
theilbar, so müssen sie gleich $q_h - 1$ sein; die Zahlen (13) sind
alle $= 1$ und es würde also folgen, dass die Zahlen $i_k$ und
also auch $J$ relativ prim zu $q_h - 1$ wären, was der Voraussetzung e) widerspricht.

    Wenn wir jetzt zu der in 2. ausgesprochenen Forderung für
die Grössen $\gamma_{k,h}$ zurückkehren, so ergiebt sich zunächst, dass
die Congruenzen (9) erfüllt sein müssen, wenn $x_k = i_k$ und die
übrigen $x = 0$ gesetzt werden, also:
\[
\gamma_{k,h}\, i_k \equiv 0\ (\mathrm{mod}\ c_h),
\]
und daraus mit Berücksichtigung von (11):
\[
i_{k,h}\, \gamma_{k,h} \equiv 0\ (\mathrm{mod}\ c_{k,h}).
\]

    Weil aber $i_{k,h}$ und $c_{k,h}$ relativ prim sind, so folgt, dass $\gamma_{k,h}$
durch $c_{k,h}$ theilbar sein muss. Wir führen also ein neues System
von ganzen Zahlen $e_{k,h}$ ein durch die Gleichungen
\weq{15}{\gamma_{k,h} = c_{k,h}\, e_{k,h}}
und suchen nun für diese Zahlen die aus 2. folgenden Bedingungen. Da die Zahlen $\gamma_{k,h}$, wie aus ihrer Definition hervorgeht, nur nach dem Modul $c_h$ bestimmt sind, so ist $e_{k,h}$ nur nach
dem Modul $\delta_{k,h}$ zu bestimmen. Nehmen wir eine von den Congruenzen (9):
\[
\gamma_{1,h}x_1 + \gamma_{2,h}x_2 + \cdots + \gamma_{\nu,h}x_\nu \equiv 0\ (\mathrm{mod}\ c_h),
\]
und führen darin (15) ein, so erhält sie die Form
\weq{16}{c_{1,h}e_{1,h}x_1 + c_{2,h}e_{2,h}x_2 + \cdots + c_{\nu,h}e_{\nu,h}x_\nu \equiv 0\ (\mathrm{mod}\ c_h).}
Hierin setzen wir nun nach (11):
\[
c_{k,h} = \frac{c_h}{\delta_{k,h}} = \frac{c_h\, i_{k,h}}{i_k}\,.
\]
Demnach sind die Congruenzen (16) gleichbedeutend mit
der Forderung, dass
%CONT>>
